\subsection{ORTHOGONAL REFLECTIONS}

Assume that V is finite dimensional. Let B be a non-degenerate bilinear form on V. By Algebra, Chap. IX, §6, no. 3, Prop. 4, B is invariant under a reflection s in V if and only if the subspaces \(V_{s}^{+}\) and \(V_{s}^{-}\) of V are orthogonal with respect to B; they are then non-isotropic. Moreover, for any non-isotropic hyperplane H in V, there is a unique reflection s that preserves B and induces the identity on H; this is the symmetry with respect to H, cf.~Algebra, Chap. IX, §6, no. 3. If a is a non-zero vector orthogonal to H, we have \(B(a,a)\neq0\) and the reflection s is given by the formula

\[
s (x) = x - 2 \frac {\mathrm{B} (x , a)}{\mathrm{B} (a , a)}. a \quad \text { for   any } x \in \mathrm{V},\tag{4}
\]

by Algebra, Chap. IX, §6, no. 4, formula (6). The reflection \(s\) is also called the orthogonal reflection with respect to H.

PROPOSITION 4. Assume that V is finite dimensional. Let B be a non-degenerate symmetric bilinear form on V, X a subspace of V and \(X^{0}\) the orthogonal complement of X with respect to B; finally, let s be the orthogonal reflection with respect to a non-isotropic hyperplane H of V. The following conditions are equivalent:

\begin{enumerate}
\def\labelenumi{(\roman{enumi})}
\item
  \(X\) is stable under \(s\) .
\item
  \(\mathrm{X}^0\) is stable under \(s\) .
\item
  H contains X or \(\mathrm{X}^0\) .
\end{enumerate}

We have \(V_{s}^{+} = H\) , and by what we have said, \(V_{s}^{-}\) is the orthogonal complement \(H^{0}\) of H with respect to B. By Prop. 3, X is stable under s if and only if \(X \subset H\) or \(H^{0} \subset X\) ; but the relation \(H^{0} \subset X\) is equivalent to \(X^{0} \subset H\) by Algebra, Chap. IX, §1, no. 6, Cor. 1 to Prop. 4. This proves the equivalence of (i) and (iii); that of (ii) and (iii) follows by interchanging the roles of X and \(X^{0}\) , since \((X^{0})^{0} = X\) .

